%% NEW 2026-10-16. The anatomy of a decode step: why some work happens once and
%% some happens L times. This is the structural fact the rest of the design follows
%% from, and it was previously only implied by a table of steps.
\begin{figure}[t]
\centering
\resizebox{\textwidth}{!}{%
\begin{tikzpicture}[
  x=1cm, y=1cm, font=\scriptsize,
  once/.style={draw, rounded corners=2pt, align=left, inner sep=2mm,
               fill=black!5, font=\scriptsize},
  lay/.style={draw, rounded corners=2pt, align=left, inner sep=1.6mm,
              fill=blue!8, draw=blue!45, text width=2.45cm, font=\scriptsize},
  plane/.style={draw=blue!50, fill=blue!10, minimum width=1.7cm,
                minimum height=0.32cm, inner sep=0pt},
  lbl/.style={font=\scriptsize\itshape, text=black!60, align=left},
  hd/.style={font=\bfseries\scriptsize, anchor=west},
  arr/.style={-{Latex[length=1.8mm]}, semithick},
]

%% ============ left: a request's KV is a stack of L planes ==============
\node[hd] at (-1.1,4.9) {a host-tier request's KV};
\foreach \y in {4.20, 3.78, 3.36, 2.94} \node[plane] at (0.3,\y) {};
\node[anchor=west, font=\tiny, text=black!55] at (1.3,4.20) {layer $0$};
\node[anchor=west, font=\tiny, text=black!55] at (1.3,3.78) {layer $1$};
\node[anchor=west, font=\tiny, text=black!55] at (1.3,3.36) {$\vdots$};
\node[anchor=west, font=\tiny, text=black!55] at (1.3,2.94) {layer $L{-}1$};

\draw[decorate, decoration={brace, amplitude=1.4mm, mirror}, black!45]
     (-0.62,2.78) -- (-0.62,4.36);
\node[lbl, anchor=east, align=right] at (-0.80,3.57)
     {same blocks,\\same request,\\\textbf{every} plane};

\node[lbl, anchor=north west, text=black!70] at (-1.1,2.35)
     {\textbf{Which} blocks a request owns does not\\
      change between layers --- only \textbf{what is}\\
      \textbf{in them} does. So addressing is resolved\\
      \emph{once}; the bytes move $L$ times.};

%% ============ middle: the once-per-step half ===========================
\node[hd] at (5.4,4.9) {once per step};
\node[once, anchor=north west, text width=5.3cm] at (5.4,4.55)
     {1.~decide each request's tier, admit\\[1pt]
      2.~release device blocks for requests\\
      \phantom{2.~}whose KV has reached the host\\[1pt]
      3.~resolve addressing for all $L$ layers\\[1pt]
      4.~apply the previous step's host writes};
\node[lbl, anchor=north west] at (5.4,2.35)
     {cost grows with the \textbf{number of}\\
      \textbf{requests} --- not with depth,\\
      not with context};

%% ============ the fan-out ==============================================
\draw[arr, black!65] (11.35,3.60) -- (12.80,3.60);
\node[font=\tiny\itshape, text=black!70, anchor=south, align=center]
     at (12.08,3.74) {one table,\\reused $L\times$};

%% ============ right: the per-layer half ================================
\node[hd] at (13.05,4.9) {per layer, $L$ times};
\foreach \x/\t in {13.05/{layer $\ell$}, 15.90/{layer $\ell{+}1$}} {
  \node[lay, anchor=north west] at (\x,4.55)
       {\textbf{a.} bring this layer's\\
        \phantom{a.} bytes across the link\\[2pt]
        \textbf{b.} attend\\[2pt]
        \textbf{c.} append the new\\
        \phantom{c.} token};
}
\node[anchor=north west, font=\scriptsize, text=black!55] at (18.80,3.35) {$\cdots$};

\node[lbl, anchor=north west] at (13.05,2.35)
     {cost grows with \textbf{bytes} and with\\
      \textbf{depth}. Step (a) is the only one\\
      that touches the link, and the only\\
      one that can be overlapped};

\end{tikzpicture}}
\caption{\textbf{Why a decode step has two halves.} A request's KV is a stack of $L$
planes, one per layer. The \emph{set of blocks} it occupies is a property of the
request and is identical in every plane; only the \emph{contents} differ. So the
expensive question --- where does this request's history live --- is answered once
per step and reused by all $L$ layers, while the bytes themselves must cross the link
once per layer. The rest of the design follows from that split: per-request work is
kept off the per-layer path, and the per-layer path is reduced to a transfer that can
be hidden.}
\label{fig:step}
\end{figure}
